\chapter{Background and Related Work}
\label{ch:background}

This chapter situates the thesis among three bodies of work: long-context
retrieval for question answering, agent-memory systems for LLMs, and black-box
optimization. It closes by stating precisely what is new here.

\section{Long-context question answering and retrieval}
When the evidence for a question exceeds the model's context window, the standard
solution is retrieval-augmented generation \citep{lewis2020rag}: an external
index returns a small set of passages that are concatenated into the prompt.
Retrieval quality is therefore decisive. Dense retrievers such as Dense Passage
Retrieval \citep{karpukhin2020dpr} and the unsupervised Contriever
\citep{izacard2022contriever} embed queries and passages into a shared space; a
sparse lexical baseline, BM25 \citep{robertson2009bm25}, remains strong when the
query shares vocabulary with the target. More recent systems learn \emph{when}
and \emph{what} to retrieve --- Self-RAG \citep{asai2023selfrag} interleaves
retrieval with self-critique, and LongRAG \citep{jiang2024longrag} pairs long
retrieval units with long-context readers --- or impose hierarchy on the corpus,
as in RAPTOR's recursive summary tree \citep{sarthi2024raptor} and HiQA's
hierarchical multi-document augmentation \citep{tao2024hiqa}. In all of these the
retrieval \emph{policy} is either fixed or trained by gradient descent on
in-domain supervision; none treats the storage-and-scoring rule as a small,
task-adaptive parameter vector tuned by black-box search, which is the object of
study here.

\section{Memory systems for LLM agents}
A parallel line of work gives an agent a persistent, evolving memory rather than
a one-shot retrieval index. MemGPT \citep{packer2023memgpt} organizes memory into
tiers (core, archival, external) and lets the LLM page information in and out
like an operating system. HippoRAG \citep{gutierrez2024hipporag} builds a
knowledge graph over the corpus and retrieves with Personalized PageRank.
MemoryLLM \citep{wang2024memoryllm} expands the model's hidden state with a
self-updatable memory pool. Practitioner frameworks such as mem0
\citep{chhikara2024mem0} and LangMem \citep{langmem2024} provide production-grade
persistent memory, and StreamingLLM \citep{xiao2024streamingllm} keeps an
unbounded stream tractable with attention sinks. These systems are sophisticated,
but their storage and retrieval \emph{policies} are hand-designed; the decision
of what to write to memory and how to weight a retrieval is not itself
optimized, and certainly not re-optimized per task. This thesis fills exactly
that gap.

\section{Black-box optimization of structured objects}
Because the storage decision is discrete and non-differentiable (an event is
either written or not), gradient descent is awkward and a black-box optimizer is
natural. Evolution Strategies scale well to such settings
\citep{salimans2017es}, and the Covariance Matrix Adaptation Evolution Strategy
(CMA-ES) \citep{hansen2016cmaes} is a robust derivative-free optimizer for
low-dimensional continuous vectors --- exactly the regime of a ten-dimensional
$\thetavec$. Bayesian optimization \citep{snoek2012bayesopt} is the main
alternative for expensive black-box objectives and has been used to tune
retrieval and embedding hyper-parameters. The novelty here is not the optimizer
but the \emph{object} it is pointed at: a memory-construction policy, tuned
directly on a downstream recall objective.

\section{Structured and graph memory}
Representing memory as a graph of typed nodes (events, entities) and edges
(temporal, mention) is a recurring idea, from knowledge-graph RAG
\citep{gutierrez2024hipporag} to entity-centric reading comprehension. Our
graph-memory implementation uses NetworkX \citep{hagberg2008networkx} and adds a
learnable retrieval score that linearly combines a graph-traversal term, an
embedding-similarity term, and a recency term, with weights that are part of
$\thetavec$. A separate exploratory strand, verbal self-reflection
(Reflexion, \citealp{shinn2023reflexion}), is discussed as a negative result in
Appendix~\ref{app:tables}.

\section{Benchmarks for long-context question answering}
Stage~3 evaluates on six published benchmarks chosen to span disjoint domains:
HotpotQA for multi-hop Wikipedia reasoning \citep{yang2018hotpotqa}, QASPER for
information-seeking questions over NLP papers \citep{dasigi2021qasper}, CUAD for
expert-annotated legal-contract review \citep{hendrycks2021cuad}, NarrativeQA for
abstractive questions over full books \citep{kocisky2018narrativeqa}, FinanceBench
for questions over SEC filings \citep{islam2023financebench}, and LongMemEval for
long-term interactive dialogue memory \citep{wu2024longmemeval}. Answers are
generated by \texttt{gpt-4o-mini} and scored by an independent Claude judge,
following the LLM-as-judge methodology and its known biases
\citep{zheng2023judge}; embeddings use the distilled MiniLM sentence encoder
\citep{wang2020minilm,reimers2019sbert}.

\section{Positioning: what is new}
\label{sec:gap}
Relative to the work above, this thesis contributes (i) a single parameter
vector $\thetavec$ that governs the \emph{entire} memory-construction pipeline
--- storage, abstraction, and retrieval scoring --- rather than any one stage;
(ii) \emph{per-task} optimization of that vector by black-box search, with the
transfer behaviour (within- vs.\ across-family) measured rather than assumed;
and (iii) a corpus-cumulative evaluation on six real benchmarks with gold
relevance signals, cross-vendor LLM judging, held-out significance testing, and a
disclosed adversarial self-audit. To our knowledge, no prior system tunes a
memory-construction policy per task with a derivative-free optimizer and reports
its transfer and its honest comparison against a full-context baseline.
